\documentclass[11pt]{article}
\usepackage{amsmath,amssymb,amsthm,xcolor,geometry,hyperref,booktabs}
\geometry{margin=1in}
\newcommand{\chg}[1]{\textcolor{blue}{#1}}
\newcommand{\tr}{\operatorname{tr}}
\newcommand{\Res}{\operatorname{Res}}
\newcommand{\arctanh}{\operatorname{arctanh}}
\newcommand{\arccosh}{\operatorname{arccosh}}
\newcommand{\Pf}{\operatorname{Pf}}
\title{Challenge 3: Geodesic in AdS/BCFT of a black hole \\ \large Solution (independent derivation)}
\date{}
\begin{document}\maketitle

\section{Strategy}
In the geodesic (WKB) approximation the one-point function of a heavy scalar operator in AdS/BCFT is
\begin{equation}
\langle\mathcal O(x)\rangle\;\simeq\;\mathcal N_{\mathcal O}\,e^{-m L_{\rm ren}},
\label{eq:wkb}
\end{equation}
where $L_{\rm ren}$ is the renormalised length of the shortest geodesic connecting the boundary insertion point to the end-of-the-world (ETW) brane, ending orthogonally on the brane (the endpoint on the brane is free, so it is extremised as well). We therefore need (i) the brane profile, (ii) the shortest geodesic from the boundary point at $\tau_E=0$ to the brane, and (iii) its renormalised length.

\section{Geometry of the Euclidean black hole}
With $L_{\rm AdS}=1$, the Euclidean BTZ metric is $ds^2=f\,d\tau_E^2+dr^2/f+r^2d\phi^2$, $f=r^2-r_0^2$. The $(\tau_E,r)$ plane is a hyperbolic disc: setting
\begin{equation}
r=r_0\cosh\rho,\qquad \theta=r_0\,\tau_E\ \ (\theta\sim\theta+2\pi),
\end{equation}
one finds
\begin{equation}
ds^2=d\rho^2+\sinh^2\!\rho\,d\theta^2+r_0^2\cosh^2\!\rho\,d\phi^2 ,
\end{equation}
so the horizon $r=r_0$ is the centre $\rho=0$ of the disc, the Euclidean time $\tau_E$ is the polar angle, and the $\phi$ circle is fibred over the disc. Distances along a radial line are $d\rho=dr/\sqrt{f}$.

\section{Brane profile}
For a two-dimensional brane in three bulk dimensions the Neumann condition $K_{ab}-Kh_{ab}=-\eta h_{ab}$ gives, after taking the trace ($K=2\eta$),
\begin{equation}
K_{ab}=\eta\,h_{ab}.
\end{equation}
Rotational symmetry means the brane is a curve $(\tau_E(\ell),r(\ell))$ in the disc, parametrised by proper length $\ell$, times the $\phi$ circle; the induced metric is $d\ell^2+r^2d\phi^2$, with
\begin{equation}
f\dot\tau_E^2+\frac{\dot r^2}{f}=1 .
\label{eq:norm}
\end{equation}
The $\phi\phi$ component of the extrinsic curvature is $K_{\phi\phi}=\frac12 n^\mu\partial_\mu(r^2)=r\,n^r$, so $K_{\phi\phi}=\eta r^2$ requires $n^r=\eta r$. In the orthonormal frame $(\sqrt f\,d\tau_E,\,dr/\sqrt f)$ the unit tangent is $(\sqrt f\dot\tau_E,\dot r/\sqrt f)$ and the unit normal is $(-\dot r/\sqrt f,\sqrt f\dot\tau_E)$, i.e.\ $n^r=f\dot\tau_E$. Hence
\begin{equation}
f\,\dot\tau_E=\eta\,r,\qquad \dot r^2=f-\eta^2r^2=(1-\eta^2)r^2-r_0^2 ,
\label{eq:brane}
\end{equation}
which also solves the $\ell\ell$ component (constant geodesic curvature $\eta$ in the hyperbolic disc). The brane therefore has a turning point $\dot r=0$ at
\begin{equation}
r_c=\frac{r_0}{\sqrt{1-\eta^2}}\;>\;r_0,\qquad\text{i.e.}\qquad \cosh\rho_c=\frac{1}{\sqrt{1-\eta^2}},\quad \tanh\rho_c=\eta,\quad \rho_c=\arctanh\eta .
\end{equation}
In the hyperbolic disc a curve of constant geodesic curvature $\eta<1$ is a hypercycle: the locus of points at fixed distance $\rho_c=\arctanh\eta$ from a geodesic $\gamma$. Here $\gamma$ is the diameter $\theta=\pm\pi/2$ (i.e.\ $\tau_E=\pm\beta/4$) and the brane lies on the far side of $\gamma$ from the insertion point, its closest approach to the centre being at $\theta=\pi$, i.e.\ $\tau_E=\beta/2$.

Two further facts confirm that this is the geometry described in the problem. First, the region kept for $\eta>0$ is the ``large'' side of the hypercycle, which contains the horizon $\rho=0$: in Lorentzian signature the brane is therefore behind the horizon of a one-sided black hole. Second, integrating (\ref{eq:brane}) from $r_c$ to $\infty$,
\begin{equation}
\Delta\tau_E=\int_{r_c}^\infty\frac{\eta\,r\,dr}{(r^2-r_0^2)\sqrt{(1-\eta^2)r^2-r_0^2}}
=\eta\int_0^\infty\frac{dw}{w^2+\eta^2r_0^2}=\frac{\pi}{2r_0}=\frac{\beta}{4},
\end{equation}
independently of $\eta$ (we substituted $w=\sqrt{(1-\eta^2)r^2-r_0^2}$). The brane thus meets the conformal boundary at $\tau_E=\beta/2\mp\beta/4=\pm\beta/4$: the BCFT lives on the cylinder $[-\beta/4,\beta/4]\times S^1$ with the boundary state at both ends, and $\tau_E=0$ is exactly its time-reflection symmetric midpoint, as stated in the problem.

\section{The dominant geodesic}
The insertion point sits on the boundary of the disc at $\theta=0$ (at any fixed $\phi$). Because both the geometry and the brane are $\phi$-independent, the shortest geodesic stays at constant $\phi$ and lives in the disc. Every path from the insertion point $p$ to the brane $H$ must cross the diameter $\gamma$ at some point $q$, and every point of $\gamma$ is at distance at least $\rho_c$ from $H$; hence
\begin{equation}
\operatorname{dist}(p,H)\;\ge\;\operatorname{dist}(p,\gamma)+\rho_c ,
\end{equation}
with equality along the diameter $\theta=0,\pi$ through the centre, which meets $H$ orthogonally at its point of closest approach (there $\dot r=0$, so the brane tangent is along $\theta$ while the geodesic is radial). The dominant geodesic therefore runs radially from the cutoff surface $r=r_\epsilon$ down to the horizon $r=r_0$ (the centre of the disc) and out again on the antipodal side to $r=r_c$:
\begin{equation}
L_\epsilon=\int_{r_0}^{r_\epsilon}\frac{dr}{\sqrt{r^2-r_0^2}}+\int_{r_0}^{r_c}\frac{dr}{\sqrt{r^2-r_0^2}}
=\arccosh\frac{r_\epsilon}{r_0}+\arccosh\frac{r_c}{r_0}
=\arccosh\frac{r_\epsilon}{r_0}+\arctanh\eta .
\end{equation}

\section{Renormalisation and result}
Near the boundary $ds^2\simeq r^2d\tau_E^2+dr^2/r^2$, so $z=1/r$ is the Fefferman--Graham coordinate and the cutoff is $r_\epsilon=1/\epsilon$. Using $\arccosh(r_\epsilon/r_0)=\ln(2r_\epsilon/r_0)+O(r_\epsilon^{-2})$ and subtracting the universal divergence $\ln(1/\epsilon)=\ln r_\epsilon$,
\begin{equation}
L_{\rm ren}=\ln\frac{2}{r_0}+\arctanh\eta=\ln\frac{2}{r_0}+\frac12\ln\frac{1+\eta}{1-\eta}.
\end{equation}
Inserting into (\ref{eq:wkb}),
\begin{equation}
\boxed{\;\langle\mathcal O(x)\rangle\;\simeq\;\mathcal N_{\mathcal O}\Big(\frac{r_0}{2}\Big)^{m}e^{-m\arctanh\eta}
=\mathcal N_{\mathcal O}\Big(\frac{r_0}{2}\Big)^{m}\Big(\frac{1-\eta}{1+\eta}\Big)^{m/2}
=\mathcal N_{\mathcal O}\Big(\frac{\pi}{\beta}\Big)^{m}\Big(\frac{1-\eta}{1+\eta}\Big)^{m/2}\;}
\end{equation}
The result is independent of the position $x$ (i.e.\ of $\phi$) by rotational symmetry. Its dependence on the data is: a power law $r_0^{m}$ in the horizon radius (equivalently $T^{m}$ in the temperature, $T=r_0/2\pi$), and a suppression $\big(\frac{1-\eta}{1+\eta}\big)^{m/2}=e^{-m\arctanh\eta}$ in the tension, which decreases monotonically from $1$ at $\eta\to0$ to $0$ at $\eta\to1^-$ (the brane recedes to infinity). The mass enters only through the exponent, $e^{-mL_{\rm ren}}$; beyond leading WKB order one should replace $m\to\Delta=1+\sqrt{1+m^2}\simeq m$.

\section{Consistency checks}
\begin{enumerate}
\item \emph{BCFT strip map.} The BCFT lives on a strip of width $W=\beta/2$ (with $\phi$ effectively non-compact in the black-hole phase). Mapping the strip $0<y<W$ to the upper half plane by $\zeta=e^{\pi w/W}$ gives $\langle\mathcal O(y)\rangle=A_{\mathcal O}\big(\frac{\pi}{2W\sin(\pi y/W)}\big)^{\Delta}$, which at the midpoint $y=W/2$ equals $A_{\mathcal O}\,(\pi/2W)^{\Delta}=A_{\mathcal O}\,(r_0/2)^{\Delta}$. This reproduces the factor $(r_0/2)^{m}$ exactly, including the $2^{-m}$.
\item \emph{Boundary one-point coefficient.} In Poincar\'e AdS$_3$ with the same brane (a hypercycle at distance $\arctanh\eta$ from the vertical plane $y=0$), the renormalised distance from $(y_0,\epsilon)$ to the brane is $\ln(2y_0)+\arctanh\eta$, giving $\langle\mathcal O\rangle=(2y_0)^{-\Delta}\big(\frac{1-\eta}{1+\eta}\big)^{\Delta/2}$; thus $A_{\mathcal O}=\big(\frac{1-\eta}{1+\eta}\big)^{\Delta/2}$ in the geodesic approximation, consistent with the tension dependence found above.
\item For $\eta\to0$ the brane becomes the geodesic diameter $\tau_E=\pm\beta/4$ and $\langle\mathcal O\rangle\to\mathcal N_{\mathcal O}(r_0/2)^{m}$.
\end{enumerate}

\section{Comparison with the provided solution}
The provided solution follows the same route (brane turning point $r_c=r_0/\sqrt{1-\eta^2}$, geodesic through the Euclidean horizon, $L_{\rm ren}=\ln(2/r_0)+\arctanh\eta$) and reaches the same result. I found no errors. \chg{Two remarks are added here rather than corrections: (i) the claim that the radial geodesic through the horizon is the \emph{dominant} one is justified above by the hypercycle distance inequality, and the geometric consistency of the setup (brane endpoints at $\tau_E=\pm\beta/4$ for all $\eta$, region kept containing the horizon for $\eta>0$) is verified explicitly; (ii) the factor $2^{-m}$ is scheme independent once the standard subtraction $\ln(1/\epsilon)$ is used, as confirmed by the strip map, and the overall normalisation $\mathcal N_{\mathcal O}$ (set to $1$ in the provided final answer) is not fixed by the geodesic approximation.}
\end{document}